%------------------------------------------------------------------------------%
\section{Fórmulas de redução ao primeiro quadrante}

Ao trabalharmos com Trigonometria e deparamo-nos com um ângulo que não se
encontra no primeiro quadrante, sempre podemos reduzi-lo de forma a encontrar o
ângulo correspondente a esse que esteja justamente no 1\textsuperscript{o} quadrante.
Usamos esse procedimento em virtude ds simetria encontrada no ciclo
trigonométrico. Mas precisamos nos atentar para o que ocorre com os sinais das
funções trigonométricas em cada quadrante. A Tabela~\ref{sinal1} e a
Figura~\ref{Sinai1} mostram os sinais em cada quadrante.

%------------------------------------------------------------------------------%
\begin{table}
  \begin{ceqn}
  \[
    \begin{array}{l<{\quad}>{\quad}c<{\quad}>{\quad}c<{\quad}>{\quad}c<{\quad}>{\quad}r}
      \hline
      &
      \text{1\textsuperscript{o} Q} &
      \text{2\textsuperscript{o} Q} &
      \text{3\textsuperscript{o} Q} &
      \text{4\textsuperscript{o} Q} \\ \hline
      \sin(x) & + & + & - & - \\
      \cos(x) & + & - & - & + \\
      \tan(x) & + & - & + & - \\ \hline
    \end{array}
  \]
  \end{ceqn}
  \caption{Sinais em cada quadrante.}
  \label{sinal1}
\end{table}
%------------------------------------------------------------------------------%

%------------------------------------------------------------------------------%
\begin{figure}
    \centering
    \includegraphics{sinais}
    \caption{Sinais em cada quadrante.}
    \label{Sinai1}
\end{figure}
%------------------------------------------------------------------------------%

%------------------------------------------------------------------------------%
\begin{figure}
  \centering
  \includegraphics[scale=2]{ciclo2p1}
  \caption{Redução do 2\textsuperscript{o} quadrante para o
    1\textsuperscript{o} quadrante.}
  \label{ciclo2p1}
\end{figure}
%------------------------------------------------------------------------------%

%------------------------------------------------------------------------------%
\begin{figure}
  \centering
  \includegraphics[scale=2]{terp1}
  \caption{Redução do 3\textsuperscript{o} quadrante para o
    1\textsuperscript{o} quadrante.}
  \label{ciclo2p2}
\end{figure}
%------------------------------------------------------------------------------%

%------------------------------------------------------------------------------%
\begin{figure}
  \centering
  \includegraphics[scale=2]{qp1}
  \caption{Redução do 4\textsuperscript{o} quadrante para o
    1\textsuperscript{o} quadrante.}
  \label{ciclo2p3}
\end{figure}
%------------------------------------------------------------------------------%

Se o ângulo com o qual estamos trabalhando for $y$ e ele estiver no
segundo quadrante, seu correspondente no 1\textsuperscript{o} quadrante será o
ângulo $x$ tal que
\[
  \pi - x = y
  \quad\text{ ou }\quad
  \ang{180} - x = y.
\]
Você não deve
se ater as fórmulas, a Figura~\ref{ciclo2p1} mostra que você pode fazer
essa redução usando o círculo trigonométrico.

Se o ângulo $y$ pertencer ao terceiro quadrante, seu correspondente
$x$ no primeiro quadrante será dado por
\[
  x + \pi = y
  \quad\text{ ou }\quad
  \ang{180} + x = y.
\]
A Figura~\ref{ciclo2p2} mostra que você pode fazer essa redução
usando o círculo trigonométrico.

Se o ângulo com o qual estamos trabalhando for $y$ e ele estiver no
quarto quadrante, seu correspondente no 1\textsuperscript{o} quadrante será o
ângulo $x$ tal que
\[
  2\pi - x = y
  \quad\text{ ou }\quad
  3\ang{60} - x = y.
\]
A Figura~\ref{ciclo2p3} mostra que você pode fazer essa redução usando o
círculo trigonométrico.

%------------------------------------------------------------------------------%
\begin{observation}
  {\color{red}
  Tente $n$
  }
\end{observation}
%------------------------------------------------------------------------------%

%------------------------------------------------------------------------------%
\begin{figure}
  \centering
  \includegraphics[scale=2]{135}
  \caption{Redução do \ang{135} ao 1\textsuperscript{o} quadrante.}
  \label{fig:reducao_135}
\end{figure}
%------------------------------------------------------------------------------%

%------------------------------------------------------------------------------%
\begin{example}
  Reduza o ângulo \ang{135}
  e calcule os valores de
  $\sin(x)$, $\cos(x)$ e $\tan(x)$ para esse ângulo.

  \solution
  O arco $x=\ang{135}$ está no $2^{\circ}$ quadrante, então para
  reduzirmos ao primeiro quadrante basta usarmos a fórmula
  \[
    \ang{180} - x = y.
  \]
  Logo,
  \[
    y
    = \ang{180} - \ang{135}
    = \ang{45}.
  \]

  Para encontrarmos as relações trigonométricas para esse ângulo, devemos
  nos atentar para o sinal das funções trigonométricas no quadrante a qual
  ele pertence. Nesse caso, temos que $x=\ang{135}$ está no segundo
  quadrante. Assim,
  \begin{gather*}
    \sin(\ang{135}) =  \sin(\ang{45}) =  \frac{\sqrt{2}}{2} \\[2mm]
    \cos(\ang{135}) = -\cos(\ang{45}) = -\frac{\sqrt{2}}{2} \\[2mm]
    \tan(\ang{135}) = -\tan(\ang{45}) = -1
  \end{gather*}
\end{example}
%------------------------------------------------------------------------------%

%------------------------------------------------------------------------------%
\begin{example}
  Reduza o ângulo
  \(\frac{11 \pi}{6}\) rad
  e calcule os valores de
  $\sin(x)$, $\cos(x)$ e $\tan(x)$ para esse ângulo.

  \solution
  Sabendo que
  \[
    \pi \text{ rad} = \ang{180},
  \]
  então
  \[
    \frac{11\pi}{6} \text{ rad}
    = \frac{11 \cdot 180}{6}
    = \ang{330}.
  \]
  Assim,
  \[
    x = \frac{11\pi}{6} \text{ rad}
  \]
  está no 4\textsuperscript{o} quadrante.
  Para reduzirmos ao primeiro quadrante, basta
  usarmos a fórmula
  \[
    y =\ang{360} - x.
  \]
  Logo,
  \[
    y
    = \ang{360} - \ang{330}
    = \ang{30}.
  \]
  Observando o sinal de cada função trigonométrica no
  4\textsuperscript{o} temos
  \begin{gather*}
    \sin(\ang{330}) = -\sin(\ang{30}) = -\frac{1}{2}          \\[2mm]
    \cos(\ang{330}) =  \cos(\ang{30}) =  \frac{\sqrt{3}}{2}   \\[2mm]
    \tan(\ang{330}) = -\tan(\ang{30}) = -\frac{\sqrt{3}}{3}
  \end{gather*}
\end{example}
%------------------------------------------------------------------------------%

%------------------------------------------------------------------------------%
\begin{example}
  Reduza o ângulo
  \(\frac{4\pi}{3}\) rad
  e calcule os valores de
  $\sin(x)$, $\cos(x)$ e $\tan(x)$ para esse ângulo.

  \solution
  Sabendo que
  \[
    x
    = \frac{4 \cdot 180}{3}
    = \ang{240},
  \]
  então $x$ está no 3\textsuperscript{o}
  quadrante e para reduzirmos ao terceiro quadrante usamos a relação
  \[
    y = x-\ang{180}.
  \]
  Logo,
  \[
    y
    = \ang{240} - \ang{180}
    = \ang{60}.
  \]
  Assim, observando o sinal de cada função trigonométrica no
  3\textsuperscript{o} temos
  \begin{gather*}
    \sin(\ang{240}) = -\sin(\ang{60}) = -\frac{\sqrt{3}}{2} \\[2mm]
    \cos(\ang{240}) =  \cos(\ang{60}) =  \frac{1}{2}        \\[2mm]
    \tan(\ang{240}) = -\tan(\ang{60}) =  \sqrt{3}
  \end{gather*}
\end{example}
%------------------------------------------------------------------------------%

